% Appendix C (two-cause rewrite, 2026-11): extended results. Every table is generated.
\chapter{Extended Results}
\label{app:extended-results}

This appendix gives the values behind the summaries of Chapter~\ref{ch:results}: the
results of every run, the retention matrices, the slices of the softmax on CIL-FUNSD and
the design choices of SBR-Recall.

\section{Results of Every Run}
\label{app:per-run}

Table~\ref{tab:perseed} lists the \gls{aa} of every run behind the main cells on CIL-CORD,
Mixed and DIL, with the mean training time of a run. Table~\ref{tab:newstreams-runs} does
the same for CIL-WildReceipt and DIL-XFUND.

\begin{table}[!htbp]
  \centering
  \caption[AA of every run on CIL-CORD, Mixed and DIL]{\textbf{AA of every run} on
  CIL-CORD, Mixed and DIL (columns: the three runs; h/run: mean training time in hours).
  Source: \emph{clean} is the protocol of Section~\ref{sec:e-protocol}; the rows marked
  $^{\dagger}$ follow an earlier protocol and belong to Appendix~\ref{app:earlier}.}
  \label{tab:perseed}
  \footnotesize% AUTO-GENERATED by scripts/rca/make_clean_tables.py -- do not edit.
\begin{tabular}{llcccc}
\toprule
Row & Source & 42 & 7 & 123 & h/run \\
\midrule
\multicolumn{6}{l}{\textit{CIL-CORD}} \\
Naive & clean & 18.8 & 18.3 & 17.8 & 1.3 \\
SBR, no memory & clean & 21.2 & 19.6 & 24.0 & 1.2 \\
CoLaR & clean & 33.6 & 32.4 & 34.4 & 1.6 \\
CoLaR-SBR & clean & 87.8 & 91.3 & 86.2 & 2.1 \\
SBR-Recall & clean & 91.4 & 92.6 & 91.9 & 1.6 \\
Joint training & clean & 92.0 & 92.0 & 93.6 & 1.6 \\
ER & clean & 27.7 & 21.1 & 23.2 & 2.7 \\
CoLaR-SBR + LexSlot & gate$^{\dagger}$ & 90.2 & 89.7 & 89.2 & 2.3 \\
CoLaR-SBR (BROS) & clean & 88.4 & 87.7 & 86.4 & 1.1 \\
LwF & clean & 18.6 & 19.0 & 17.8 & 1.7 \\
DER++ & clean & 22.8 & 21.3 & 19.5 & 6.2 \\
CODA-Prompt & clean & 7.5 & 2.5 & 11.9 & 2.5 \\
CL-LoRA (2025) & clean & 18.9 & 17.7 & 18.7 & 3.4 \\
CPFD (2023) & clean & 18.5 & 0.0 & 0.3 & 0.9 \\
IS3 (2024) & clean & 5.7 & 3.5 & 6.4 & 2.2 \\
LCA (2026) & clean & 7.9 & 10.1 & 8.3 & 2.0 \\
\midrule
\multicolumn{6}{l}{\textit{Mixed}} \\
Naive & clean & 42.5 & 42.6 & 41.3 & 1.3 \\
CoLaR & clean & 67.1 & 77.4 & 79.8 & 1.9 \\
CoLaR-SBR & clean & 88.5 & 89.2 & 89.7 & 1.8 \\
Joint training & clean & 88.4 & 86.9 & 87.4 & 0.8 \\
ER & clean & 72.5 & 78.5 & 78.0 & 3.3 \\
CoLaR-SBR + LexSlot & gate$^{\dagger}$ & 88.5 & 90.0 & 90.3 & 2.5 \\
SBR-Recall & clean & 85.1 & 83.6 & 85.5 & 4.9 \\
LwF & clean & 43.6 & 45.8 & 43.2 & 1.7 \\
DER++ & clean & 68.4 & 75.8 & 76.8 & 4.9 \\
CODA-Prompt & clean & 8.4 & 7.8 & 9.4 & 1.9 \\
CL-LoRA (2025) & clean & 41.5 & 40.8 & 41.5 & 3.0 \\
CPFD (2023) & clean & 0.0 & 0.1 & 6.0 & 1.0 \\
IS3 (2024) & clean & 0.0 & 0.0 & 42.1 & 1.7 \\
LCA (2026) & clean & 26.4 & 40.3 & 35.1 & 1.8 \\
\midrule
\multicolumn{6}{l}{\textit{DIL}} \\
Naive & clean & 39.6 & 39.0 & 39.0 & 0.8 \\
CoLaR & clean & 86.5 & 87.2 & 85.9 & 1.3 \\
CoLaR-SBR & clean & 85.7 & 85.5 & 84.9 & 1.4 \\
Joint training & clean & 86.5 & 86.7 & 87.0 & 0.9 \\
ER & clean & 86.7 & 87.2 & 87.0 & 2.1 \\
CoLaR-SBR + LexSlot & gate$^{\dagger}$ & 87.2 & 87.2 & 86.6 & 1.4 \\
SBR-Recall & clean & 32.1 & 33.4 & 30.1 & 0.6 \\
LwF & clean & 38.7 & 36.5 & 37.7 & 1.3 \\
DER++ & clean & 89.4 & 86.1 & 89.1 & 4.2 \\
CODA-Prompt & clean & 31.1 & 31.1 & 30.8 & 1.5 \\
CL-LoRA (2025) & clean & 43.4 & 42.4 & 41.8 & 1.9 \\
CPFD (2023) & clean & 0.0 & 0.0 & 0.0 & 0.6 \\
IS3 (2024) & clean & 38.7 & 38.5 & 38.1 & 1.0 \\
LCA (2026) & clean & 43.2 & 41.7 & 38.2 & 0.9 \\
\bottomrule
\end{tabular}

\end{table}

\begin{table}[!htbp]
  \centering
  \caption[Results of every run on CIL-WildReceipt and DIL-XFUND]{\textbf{Results of every
  run} on CIL-WildReceipt and DIL-XFUND.}
  \label{tab:newstreams-runs}
  \footnotesize% AUTO-GENERATED by scripts/rca/make_clean_tables.py -- do not edit.
\begin{tabular}{@{}llcccccc@{}}
\toprule
 & & \multicolumn{3}{c}{AA, seed} & \multicolumn{3}{c}{Full-label F1, seed} \\
\cmidrule(lr){3-5}\cmidrule(l){6-8}
Stream & Method & 42 & 7 & 123 & 42 & 7 & 123 \\
\midrule
CIL-WildReceipt & CoLaR-SBR & 79.3 & 79.7 & 79.2 & 79.3 & 79.8 & 79.3 \\
CIL-WildReceipt & CoLaR, shared softmax & 28.3 & 19.5 & 20.6 & 34.4 & 30.8 & 32.4 \\
CIL-WildReceipt & SBR-Recall & 80.9 & 81.5 & 80.7 & 81.1 & 81.5 & 80.6 \\
CIL-WildReceipt & ER & 20.6 & 19.9 & 20.4 & 30.3 & 30.8 & 31.3 \\
CIL-WildReceipt & DER++ (250 documents) & 18.9 & 19.9 & 20.4 & 28.2 & 31.2 & 31.7 \\
CIL-WildReceipt & LwF & 19.6 & 19.5 & 20.1 & 31.3 & 30.6 & 31.3 \\
CIL-WildReceipt & CL-LoRA & 18.0 & 18.0 & 18.7 & 29.5 & 29.6 & 30.4 \\
CIL-WildReceipt & Naive fine-tuning & 19.4 & 19.6 & 19.6 & 30.0 & 31.2 & 31.3 \\
CIL-WildReceipt & Joint training (reference) & 82.6 & 84.1 & 83.8 & 82.6 & 84.5 & 84.0 \\
\midrule
DIL-XFUND & CoLaR-SBR & 76.2 & 76.2 & 77.0 & n/a & n/a & n/a \\
DIL-XFUND & CoLaR, shared softmax & 75.0 & 77.8 & 77.8 & n/a & n/a & n/a \\
DIL-XFUND & ER & 79.9 & 80.0 & 78.2 & n/a & n/a & n/a \\
DIL-XFUND & DER++ (250 documents) & 80.2 & 80.5 & 79.8 & n/a & n/a & n/a \\
DIL-XFUND & LwF & 71.9 & 72.2 & 71.9 & n/a & n/a & n/a \\
DIL-XFUND & CL-LoRA & 57.5 & 55.6 & 47.7 & n/a & n/a & n/a \\
DIL-XFUND & Naive fine-tuning & 70.8 & 70.8 & 72.0 & n/a & n/a & n/a \\
DIL-XFUND & Joint training (reference) & 79.8 & 79.1 & 79.7 & n/a & n/a & n/a \\
\bottomrule
\end{tabular}

\end{table}

\section{Retention Matrices}
\label{app:matrices}

Tables~\ref{tab:mat-cord} and~\ref{tab:mat-mixed} give the retention matrices behind
Figure~\ref{fig:heatmaps} and behind the Mixed results (mean over the three runs; row $i$:
after training session $i$; column $j$: session $j$). On CIL-CORD, the diagonal of CoLaR
with a shared softmax falls from session~3 on, whereas CoLaR-SBR and SBR-Recall learn
every session and keep it. On Mixed without a memory, SBR-Recall loses the first two
sessions as soon as the schema changes and does not recover them when FUNSD returns in
session~6.

\begin{table}[!htbp]
  \centering
  \caption[Retention matrices on CIL-CORD]{\textbf{Retention matrices on CIL-CORD}
  (session-restricted F1).}
  \label{tab:mat-cord}
  \footnotesize
  % AUTO-GENERATED by scripts/rca/make_clean_tables.py -- do not edit.
\begin{tabular}{l*{5}{c}}
\multicolumn{6}{l}{\textbf{CoLaR, shared softmax, CIL-CORD} (seeds 42/7/123, clean)} \\
\toprule
$i\backslash j$ & 1 & 2 & 3 & 4 & 5 \\
\midrule
1 & 93.9 &  &  &  &  \\
2 & 72.3 & 70.8 &  &  &  \\
3 & 42.2 & 73.6 & 25.8 &  &  \\
4 & 29.7 & 69.3 & 57.2 & 8.4 &  \\
5 & 20.9 & 59.2 & 43.3 & 34.8 & 9.2 \\
\bottomrule
\end{tabular}
\\[8pt]
  % AUTO-GENERATED by scripts/rca/make_clean_tables.py -- do not edit.
\begin{tabular}{l*{5}{c}}
\multicolumn{6}{l}{\textbf{CoLaR-SBR, CIL-CORD} (seeds 42/7/123, clean)} \\
\toprule
$i\backslash j$ & 1 & 2 & 3 & 4 & 5 \\
\midrule
1 & 93.9 &  &  &  &  \\
2 & 93.7 & 92.3 &  &  &  \\
3 & 91.7 & 88.4 & 89.5 &  &  \\
4 & 91.0 & 86.8 & 90.1 & 89.3 &  \\
5 & 90.8 & 83.4 & 87.9 & 86.2 & 93.8 \\
\bottomrule
\end{tabular}
\\[8pt]
  % AUTO-GENERATED by scripts/rca/make_clean_tables.py -- do not edit.
\begin{tabular}{l*{5}{c}}
\multicolumn{6}{l}{\textbf{SBR-Recall, CIL-CORD} (seeds 42/7/123, clean)} \\
\toprule
$i\backslash j$ & 1 & 2 & 3 & 4 & 5 \\
\midrule
1 & 93.9 &  &  &  &  \\
2 & 94.1 & 88.4 &  &  &  \\
3 & 94.1 & 89.2 & 92.1 &  &  \\
4 & 94.1 & 89.7 & 92.3 & 89.3 &  \\
5 & 94.1 & 89.7 & 92.3 & 89.1 & 94.5 \\
\bottomrule
\end{tabular}

\end{table}

\begin{table}[!htbp]
  \centering
  \caption[Retention matrices on Mixed]{\textbf{Retention matrices on Mixed}
  (session-restricted F1).}
  \label{tab:mat-mixed}
  \footnotesize
  % AUTO-GENERATED by scripts/rca/make_clean_tables.py -- do not edit.
\begin{tabular}{l*{6}{c}}
\multicolumn{7}{l}{\textbf{CoLaR-SBR, Mixed} (seeds 42/7/123, clean)} \\
\toprule
$i\backslash j$ & 1 & 2 & 3 & 4 & 5 & 6 \\
\midrule
1 & 84.6 &  &  &  &  &  \\
2 & 87.1 & 89.9 &  &  &  &  \\
3 & 86.7 & 89.4 & 81.4 &  &  &  \\
4 & 86.4 & 87.5 & 78.9 & 95.9 &  &  \\
5 & 85.1 & 89.0 & 77.8 & 95.1 & 94.0 &  \\
6 & 87.1 & 90.5 & 79.2 & 95.3 & 94.2 & 88.5 \\
\bottomrule
\end{tabular}
\\[8pt]
  % AUTO-GENERATED by scripts/rca/make_clean_tables.py -- do not edit.
\begin{tabular}{l*{6}{c}}
\multicolumn{7}{l}{\textbf{SBR-Recall, Mixed, 50-document memory} (seeds 42/7/123, clean)} \\
\toprule
$i\backslash j$ & 1 & 2 & 3 & 4 & 5 & 6 \\
\midrule
1 & 85.4 &  &  &  &  &  \\
2 & 85.8 & 88.3 &  &  &  &  \\
3 & 65.8 & 61.0 & 82.2 &  &  &  \\
4 & 69.0 & 49.8 & 77.9 & 95.7 &  &  \\
5 & 70.0 & 51.9 & 78.8 & 96.0 & 94.9 &  \\
6 & 81.7 & 76.0 & 80.0 & 95.9 & 95.4 & 79.5 \\
\bottomrule
\end{tabular}
\\[8pt]
  % AUTO-GENERATED by scripts/rca/make_clean_tables.py -- do not edit.
\begin{tabular}{l*{6}{c}}
\multicolumn{7}{l}{\textbf{SBR-Recall, Mixed, no memory} (seeds 42/7/123, clean)} \\
\toprule
$i\backslash j$ & 1 & 2 & 3 & 4 & 5 & 6 \\
\midrule
1 & 85.4 &  &  &  &  &  \\
2 & 85.3 & 89.9 &  &  &  &  \\
3 & 0.0 & 0.0 & 81.6 &  &  &  \\
4 & 0.0 & 0.0 & 22.4 & 96.0 &  &  \\
5 & 0.0 & 0.0 & 11.4 & 96.2 & 95.7 &  \\
6 & 0.0 & 0.0 & 11.4 & 96.0 & 95.2 & 0.0 \\
\bottomrule
\end{tabular}

\end{table}

\section{Slicing the Softmax on CIL-FUNSD}
\label{app:funsd}

On CIL-FUNSD, the question, answer and header types of FUNSD arrive in two sessions, and
the forms of session~2 also contain questions, which are labelled ``O''.
Table~\ref{tab:funsd} reports the repairs of Section~\ref{sec:r-h2} on this stream
together with two slices of the softmax. The replay-side slice keeps, for a stored
document, all classes that were known when it was stored. The current-side slice removes
the logits of the old classes only for the tokens labelled ``O''; the tokens with an
entity label still use the full softmax. Combining both slices is a partial-label control
inspired by SS-IL~\citep{ahn2021ssil} and ER-ACE~\citep{caccia2022erace}, and not a full
reproduction of them. As shown in Table~\ref{tab:funsd}, CoLaR-SBR is level with joint
training, whereas the memory relabelled with the ground-truth labels stays at the level
of the shared softmax, even with fifty stored forms. In contrast, slicing both losses
closes most of the gap on this stream, and slicing only the current loss closes a part of
it.

\begin{table}[!htbp]
  \centering
  \caption[CIL-FUNSD]{\textbf{CIL-FUNSD} ($10$ stored forms per session; mean $\pm$ s.d.).
  $R[2,1]$: F1 of session~1 after session~2 (range over the runs). $^{r}$Uses the
  ground-truth labels of the stored forms.}
  \label{tab:funsd}
  \small% AUTO-GENERATED by scripts/rca/make_clean_tables.py -- do not edit.
\begin{tabular}{@{}lccc@{}}
\toprule
Method & Full-label F1 & AA & $R[2,1]$ \\
\midrule
Naive fine-tuning & $50.8{\pm}0.2$ & $43.9{\pm}0.5$ & 0.0--0.0 \\
CoLaR, shared softmax & $52.2{\pm}0.6$ & $45.0{\pm}0.2$ & 0.0--1.0 \\
CoLaR, true labels before each session$^{r}$ & $52.2{\pm}0.2$ & $44.5{\pm}0.2$ & 0.0--0.0 \\
CoLaR-SBR & $\underline{87.2{\pm}0.6}$ & $\underline{87.4{\pm}0.6}$ & 85.0--86.6 \\
CoLaR-SBR, constant background logit & $87.2{\pm}0.3$ & $\mathbf{87.7{\pm}0.3}$ & 85.5--86.0 \\
Joint training (reference) & $87.5{\pm}0.8$ & $87.9{\pm}0.7$ & -- \\
\midrule
true labels before each session, 50 forms$^{r}$ & $54.8{\pm}6.4$ & $49.8{\pm}10.1$ & 0.0--34.6 \\
CoLaR + MiB-style unbiased CE & $\mathbf{87.3{\pm}0.4}$ & $\mathbf{87.7{\pm}0.0}$ & 84.0--87.7 \\
CoLaR, replay and current losses sliced & $83.7{\pm}3.0$ & $84.4{\pm}2.8$ & 74.9--83.8 \\
SBR, no memory & $77.6{\pm}3.8$ & $79.7{\pm}3.4$ & 65.0--75.2 \\
CoLaR, current loss sliced & $78.0{\pm}6.1$ & $79.1{\pm}5.5$ & 56.3--79.7 \\
\bottomrule
\end{tabular}

\end{table}

\section{Design Choices of SBR-Recall}
\label{app:recall-design}

Table~\ref{tab:ablext} gathers two design choices of SBR-Recall. \emph{Capacity of the
updates.} An update of rank $8$ on the query, key and value maps learns its weakest
session only to the value in the column min $R[j,j]$, and adding the memory does not
repair it. In contrast, at rank $64$ with the attention output and the feed-forward maps,
every session is learned. Hence, the capacity of the update was the limit, and we use rank
$64$; this is the one setting that we chose on CIL-CORD. \emph{Across schemas.} Without a
memory on Mixed, an old block is trained towards its background on every token of the
documents of another schema, and nothing trains it towards its own tags; hence, it stops
claiming its entities and does not recover when its dataset returns: both minima of this
row are zero. Since the encoder is frozen, this is an effect of the readout and not of
drift. In contrast, with the $50$-document memory, SBR-Recall reaches $\nRecallMixMem$
\gls{aa} against $\nRecallMixNoMem$ without it. Thus, the updates do not replace the
memory across schemas.

\begin{table}[!htbp]
  \centering
  \caption[Design choices of SBR-Recall]{\textbf{Design choices of SBR-Recall} (AA, mean
  $\pm$ s.d.). Min $R[j,j]$: the lowest F1 of a session right after it is learned; min
  $R[T,j]$: the lowest final F1 of a session, both over all runs.}
  \label{tab:ablext}
  \small% AUTO-GENERATED by scripts/rca/make_clean_tables.py -- do not edit.
\begin{tabular}{@{}lccc@{}}
\toprule
Variant & AA & min $R[j,j]$ & min $R[T,j]$ \\
\midrule
\multicolumn{4}{@{}l}{\emph{SBR-Recall delta capacity, CIL-CORD}} \\
\quad rank 8, Q/K/V, with memory & $81.3{\pm}2.5$ & 45.3 & 59.4 \\
\quad rank 8, Q/K/V, no memory & $77.5{\pm}8.6$ & 32.9 & 34.3 \\
\quad rank 64, Q/K/V + dense, no memory & $91.9{\pm}0.6$ & 85.6 & 85.6 \\
\midrule
\multicolumn{4}{@{}l}{\emph{SBR-Recall across schemas, Mixed (rank 64)}} \\
\quad no memory & $33.8{\pm}2.1$ & 0.0 & 0.0 \\
\quad $50$-document memory & $84.8{\pm}1.0$ & 77.2 & 75.1 \\
\bottomrule
\end{tabular}

\end{table}
